\section{A produced model of the entire original relative embedding}

\begin{lemma}[Arithmetic descent starts from the original relative geometry]
\label{lem:kn-relative-arithmetic-original-projective-stage}
\lean{EnvelopeProperty.Vanishing.RelativeArithmeticOriginalProjectiveModel.exists_original_projective_stage}
\uses{lem:kn-relative-arithmetic-affine-base-model,lem:kn-relative-arithmetic-fixed-projective-ambient,lem:kn-relative-arithmetic-smooth-closed-model}
\leanok
Let $k$ have characteristic zero, let $R$ be an arbitrary finite-type
$k$-algebra, and suppose
$f:X\to\operatorname{Spec}k$ is separated and smooth of relative
dimension $n$.  Let $g:X\to\operatorname{Spec}R$ be an original
projective morphism compatible with $f$, and fix an original closed
embedding $j:X\hookrightarrow\mathbb P_R^d$ over $g$.
Then the original affine base internally produces a finite arithmetic
algebra $B/A$, and the original whole closed embedding internally
produces a finite coefficient stage and a smooth relative-dimension-$n$
model of $X$ closed in the corresponding scalar extension of
$\mathbb P_B^d$.  The entire original $j$ and all structure morphisms
are recovered under the original field fibre.
\end{lemma}
\begin{proof}
Apply the affine-base model lemma to the given $R$ to obtain the
actual coefficient ring $A\subset k$, the separate finitely
presented algebra $B$, and the whole original tensor isomorphism.
The prescribed-ambient lemma makes $\mathbb P_B^d$ a
finite-presentation model of the entire original $\mathbb P_R^d$.
The original $f$ is quasi-compact: it is the composite of the
projective $g$ with the affine morphism
$\operatorname{Spec}R\to\operatorname{Spec}k$.  It is
quasi-separated because it is separated, and it is locally of
finite presentation because it is smooth.  Thus the smooth closed
model lemma applies to this actual ambient model and the given
closed immersion $j$.  It produces the stated coefficient stage,
whole smooth model, and closed embedding, including the original
recovery equation.  None of these steps uses properness over $k$;
projectivity over the original $R$ supplies only the required
quasi-compactness.
\end{proof}
